\subsection{Executed cover: from dispatch to inference}\label{sec:cover-study}
An honest tool call follows a model decision. Controller-inserted cover can skip that wait. We test the resulting scheduling cue on actual encrypted records, then compare live delays with cover generated by another planner. Timing-aware dummy design is established in WF~\cite{shmatikov2006adaptive,juarez2016wtfpad,gong2020front}; here its execution and resource consequences are measured together.

\paragraph{Design and observer}
Qwen3:8B and Llama3.1:8B each run the same 300 synthetic ticket fixtures in five arms: no cover, immediate cover, historical-gap dispatch that blocks the chain, the same dispatch concurrent with inference, and a complete shadow execution of the same model. The split is 120 training, 60 validation and 120 test fixtures per cell, balanced across six requested actions. LangGraph~\cite{langchain2025langgraph} chooses checks, actions and recovery against seven local MCP Streamable HTTP services over TLS~1.3/HTTP/2. The observer sees destinations, ciphertext counts and TLS-record completion times within supplied call windows. Same-service calls serialize. Calls to different services and planner inference can overlap. Every cover executes on isolated state and stays hidden from the real planner.

No-cover training runs supply the frozen historical gaps and durations. Sequential and concurrent arms wait for these draws and hold responses for sampled minimum durations. The shadow ticket independently selects one of all six actions, including the real action; either chain can start first. Size and adaptive attacks are selected on validation data. Paired tests use the 120 held-out fixtures, with separate Holm families for the planned questions (Supplement S12).

\begin{table*}[t]\centering\small
\caption{Live timing on 120 held-out paired fixtures per model and arm (six-class reference 20/120). Size and Adapt.\ are validation-selected attacks. RF is the condition-refitted size/destination forest; GMM removes cover without role labels. Bytes includes complete unique connections, including setup/teardown. Latency is the ratio of mean real-chain admission-to-finish times; phase and cache history remain part of these descriptive costs. Goal/State counts requested-goal success and whole-state correctness. Cover is isolated and hidden from the real planner.}
\label{tab:live-timing-main}
\setlength{\tabcolsep}{5pt}
\begin{tabularx}{\textwidth}{@{}lYrrrrrrr@{}}\toprule
Model & Policy & Size & Adapt. & RF & GMM & Bytes$\times$ & Latency$\times$ & Goal/State\\\midrule
Qwen3:8B & No cover & 117 & 117 & 116 & -- & 1.00 & 1.00 & 113/114 \\
Qwen3:8B & Immediate & 75 & 112 & 59 & 112 & 1.15 & 1.23 & 111/112 \\
Qwen3:8B & History, sequential & 76 & 93 & 61 & 53 & 1.15 & 1.47 & 112/113 \\
Qwen3:8B & History, concurrent & 66 & 99 & 52 & 38 & 1.15 & 1.14 & 111/112 \\
Qwen3:8B & Planner shadow & 54 & 54 & 70 & 34 & 1.31 & 1.65 & 112/113 \\
\midrule
Llama3.1:8B & No cover & 110 & 110 & 112 & -- & 1.00 & 1.00 & 82/70 \\
Llama3.1:8B & Immediate & 80 & 108 & 80 & 108 & 1.13 & 0.88 & 83/67 \\
Llama3.1:8B & History, sequential & 82 & 104 & 82 & 14 & 1.13 & 1.15 & 84/67 \\
Llama3.1:8B & History, concurrent & 68 & 106 & 68 & 66 & 1.13 & 0.84 & 82/67 \\
Llama3.1:8B & Planner shadow & 64 & 64 & 64 & 43 & 1.35 & 1.18 & 82/67 \\
\bottomrule\end{tabularx}
\end{table*}


\paragraph{Record timing reproduces cover removal}
Independent TCP reassembly matches all 773,450 TLS records across the collection. An unsupervised mixture on training gaps supplies a removal threshold without cover-role labels; the remaining trace goes to the frozen no-cover task classifier. It restores 112/120 Qwen and 108/120 Llama predictions from size-only references of 75 and 80. Both meet the prespecified replication rule: gains exceed 24/120 and Holm-adjusted $p\leq4.3\times10^{-6}$. Recomputing gaps with packet-capture timestamps gives exactly the same two scores.

\paragraph{Historical draws are not an achieved timing distribution}
With sequential history-based dispatch, the simple mixture attack falls to 53 and 14, but validation-selected context removal or joint decoding reaches 93 and 104, against size-only 76 and 82. The paired gains are 14.2 and 18.3 points. Only Llama passes Holm correction ($p=.0559$ and $.00078$). Concurrent dispatch gives 99 and 106 against 66 and 68, gains of 27.5 and 31.7 points, both corrected-significant ($p=1.7\times10^{-5}$ and $9.7\times10^{-9}$).

The live policy implements minimum waits, not exact marginal matching. For Qwen's sequential cover, a mean sampled response duration of 40.5\,ms becomes 100.4\,ms at the record observer. Concurrent dispatch reaches 171.2\,ms. Transport and overlapping work contribute beyond the hold. The earlier receipt-level counterfactuals (Supplements S10--S11) therefore motivate the policy but do not predict its achieved cost or protection. Context attacks also use destination, position and size, so their gain is an adapted attack result, not a timing-only causal effect.

\paragraph{Planner-generated cover changes the tradeoff}
For shadow execution, both validation selectors choose the same size-only attack: 54/120 for Qwen and 64/120 for Llama. The selected comparison therefore trivially meets the frozen $\pm10$-point equivalence rule. Substantial inference remains: the prespecified size/destination RF scores 70 and 64, and the chain-separation attack scores 49 and 34. It exactly recovers the real chain in only 29 and 32 fixtures. No tested timing method exceeds the largest observed size-only score, but these results leave stronger role-separation attacks open.

Shadow execution uses 2.01$\times$ and 1.99$\times$ model tokens and 1.68$\times$ and 1.71$\times$ call-window bytes. Its mean real-chain latency is 4.80 and 3.94\,s, versus 2.91 and 3.33\,s without cover. Whole-connection byte factors in Table~\ref{tab:live-timing-main} include the fixed initialization cost. Phase order and persistent caches affect latency. Shadow and real tickets can have different cache histories. Requested-goal and whole-state outcomes remain separate columns, with no equivalence claim about business utility.

\paragraph{Execution semantics remain a separate boundary}
The original six-policy cohort isolates this question (Supplement S8). Proxy, read and hidden-isolated cover retain Qwen state integrity at 115--116/120. Deliberately adverse live writes reduce it to 7/120. Llama falls from 69 to 3. Of 102 and 76 completed live-write tasks, 95 and 73 have collateral errors. These constructed writes test the state metric, not production cover harms. The new timing study keeps cover isolated throughout: its privacy--cost comparison concerns dispatch and planner work, while the earlier control shows why business invariants need their own measurement.
